\begin{frame}{후회 137을 ``예측''하기: 손실 분해}
  $\varepsilon=0.1$ 이면 $2000 \times 0.1 = 200$ 스텝을
  무작위 팔에 사용.
  \vskip0.4em
  \begin{center}
    순수 탐험 손실 $= 200 \times 0.5 = \mathbf{100}$ \\[0.5em]
    초기 불확실성 손실 $\approx 137 - 100 = \mathbf{37}$
  \end{center}
  \begin{itemize}
    \item \textbf{탐험 손실}(100): ``활용을 포기한 대가'' ---
      무작위 팔에 쓴 스텝마다 기대 0.5 손실.
    \item \textbf{초기 불확실성 손실}(37): ``최선 팔을 \textbf{아직
      못 찾은} 초반''의 손실 --- 추정치가 불신실한 초기 구간에서
      greedy 선택도 가끔 오판.
  \end{itemize}
  \vskip0.4em
  \begin{block}{곡선의 모양}
    $\varepsilon$ 줄이면 ``탐험 손실''은 줄지만 ``초기 불확실성
    손실''은 커짐. 두 손실의 합이 \textbf{최소}가 되는 지점
    = 곡선의 꼭대기.
  \end{block}
\end{frame}
